\documentclass[aspectratio=169,xcolor=table]{beamer}

\usepackage[utf8]{inputenc}
\usepackage[T1]{fontenc}
\usepackage[french]{babel}
\usepackage{booktabs}
\usepackage{graphicx}
\usepackage{array}

\definecolor{emotycText}{HTML}{1A1A1A}
\definecolor{emotycMuted}{HTML}{71717A}
\definecolor{emotycAccent}{HTML}{3B3B3F}
\definecolor{ttkBlue}{HTML}{4F46E5}
\definecolor{cyberOrange}{HTML}{B45309}
\definecolor{softSurface}{HTML}{F7F7F8}

\setbeamertemplate{navigation symbols}{}
\setbeamertemplate{footline}[frame number]
\setbeamercolor{normal text}{fg=emotycText,bg=white}
\setbeamercolor{frametitle}{fg=emotycAccent,bg=white}
\setbeamerfont{frametitle}{series=\bfseries,size=\Large}
\setbeamerfont{title}{series=\bfseries,size=\LARGE}
\setbeamerfont{subtitle}{size=\normalsize}
\setbeamerfont{footline}{size=\scriptsize}
\setbeamersize{text margin left=8mm,text margin right=8mm}

\newcommand{\InputHeatmapTable}{%
  \IfFileExists{figures/heatmap_delta_table.tex}{%
    \begin{tabular}{lrrccc}
\toprule
\textbf{Label} & \textbf{TTK} & \textbf{Cyber} & \textbf{$\Delta F_1$} & \textbf{$\Delta P$} & \textbf{$\Delta R$} \\
\midrule
Émo & 5\,374 & 399 & \cellcolor[HTML]{C7EA97}\textcolor{black}{+0.27} & \cellcolor[HTML]{CFEB8F}\textcolor{black}{+0.31} & \cellcolor[HTML]{C0EAA1}\textcolor{black}{+0.23} \\
Base & 4\,133 & 364 & \cellcolor[HTML]{E6EC7E}\textcolor{black}{+0.37} & \cellcolor[HTML]{CBEA92}\textcolor{black}{+0.29} & \cellcolor[HTML]{EEDA6E}\textcolor{black}{+0.43} \\
Montrée & 971 & 301 & \cellcolor[HTML]{EEE172}\textcolor{black}{+0.42} & \cellcolor[HTML]{EEDF71}\textcolor{black}{+0.42} & \cellcolor[HTML]{EEE373}\textcolor{black}{+0.41} \\
Colère & 1\,180 & 290 & \cellcolor[HTML]{EFCD67}\textcolor{black}{+0.46} & \cellcolor[HTML]{EEDF71}\textcolor{black}{+0.42} & \cellcolor[HTML]{F1BB5E}\textcolor{black}{+0.49} \\
Dégoût & 52 & 80 & \cellcolor[HTML]{BBEAB5}\textcolor{black}{+0.14} & \cellcolor[HTML]{B62020}\textcolor{white}{+1.00} & \cellcolor[HTML]{C3ECC9}\textcolor{black}{+0.08} \\
Autre & 1\,270 & 62 & \cellcolor[HTML]{E4602C}\textcolor{white}{+0.76} & \cellcolor[HTML]{E25727}\textcolor{white}{+0.78} & \cellcolor[HTML]{E97B39}\textcolor{black}{+0.68} \\
Désignée & 1\,499 & 55 & \cellcolor[HTML]{EFA24F}\textcolor{black}{+0.56} & \cellcolor[HTML]{F1B55B}\textcolor{black}{+0.50} & \cellcolor[HTML]{ED9447}\textcolor{black}{+0.60} \\
Suggérée & 1\,909 & 54 & \cellcolor[HTML]{E5682F}\textcolor{black}{+0.73} & \cellcolor[HTML]{E97C3A}\textcolor{black}{+0.67} & \cellcolor[HTML]{E5642D}\textcolor{black}{+0.74} \\
Comportementale & 1\,242 & 36 & \cellcolor[HTML]{E77134}\textcolor{black}{+0.71} & \cellcolor[HTML]{EC9044}\textcolor{black}{+0.62} & \cellcolor[HTML]{E4602C}\textcolor{white}{+0.76} \\
Joie & 888 & 30 & \cellcolor[HTML]{F1BA5D}\textcolor{black}{+0.49} & \cellcolor[HTML]{F1B95D}\textcolor{black}{+0.49} & \cellcolor[HTML]{F1BB5E}\textcolor{black}{+0.49} \\
Complexe & 568 & 14 & \cellcolor[HTML]{EFA451}\textcolor{black}{+0.55} & \cellcolor[HTML]{E87637}\textcolor{black}{+0.69} & \cellcolor[HTML]{CDEB90}\textcolor{black}{+0.30} \\
Tristesse & 673 & 14 & \cellcolor[HTML]{E87737}\textcolor{black}{+0.69} & \cellcolor[HTML]{E4622D}\textcolor{black}{+0.75} & \cellcolor[HTML]{ED9347}\textcolor{black}{+0.60} \\
Jalousie & 7 & 6 & \cellcolor[HTML]{D3EEDE}\textcolor{black}{0.00} & \cellcolor[HTML]{D3EEDE}\textcolor{black}{0.00} & \cellcolor[HTML]{D3EEDE}\textcolor{black}{0.00} \\
Surprise & 824 & 5 & \cellcolor[HTML]{F0A853}\textcolor{black}{+0.54} & \cellcolor[HTML]{EE9B4B}\textcolor{black}{+0.58} & \cellcolor[HTML]{F1BB5E}\textcolor{black}{+0.49} \\
Peur & 1\,047 & 4 & \cellcolor[HTML]{EA833D}\textcolor{black}{+0.66} & \cellcolor[HTML]{E5662E}\textcolor{black}{+0.74} & \cellcolor[HTML]{EDED79}\textcolor{black}{+0.39} \\
Culpabilité & 19 & 3 & \cellcolor[HTML]{D3EEDE}\textcolor{black}{0.00} & \cellcolor[HTML]{D3EEDE}\textcolor{black}{0.00} & \cellcolor[HTML]{D3EEDE}\textcolor{black}{0.00} \\
Fierté & 202 & 3 & \cellcolor[HTML]{C9ECD1}\textcolor{black}{-0.05} & \cellcolor[HTML]{C0EBC4}\textcolor{black}{+0.09} & \cellcolor[HTML]{BFEAA3}\textcolor{black}{-0.22} \\
Embarras & 162 & 2 & \cellcolor[HTML]{E9803C}\textcolor{black}{+0.66} & \cellcolor[HTML]{E4632D}\textcolor{black}{+0.75} & \cellcolor[HTML]{BFEAA3}\textcolor{black}{-0.22} \\
Admiration & 211 & 0 & \cellcolor[HTML]{E5632D}\textcolor{black}{+0.75} & \cellcolor[HTML]{D93F20}\textcolor{white}{+0.85} & \cellcolor[HTML]{E9803C}\textcolor{black}{+0.66} \\
\midrule
\rowcolor[HTML]{3B3B3F}
\multicolumn{6}{l}{\textcolor{white}{\textbf{Métriques globales}}} \\
Macro F1 & 27\,911 & 781 & \cellcolor[HTML]{EFCC66}\textcolor{black}{+0.4570} & & \\
Micro F1 & 27\,911 & 781 & \cellcolor[HTML]{EEDC6F}\textcolor{black}{+0.4246} & & \\
\bottomrule
\end{tabular}
%
  }{%
    \begin{tabular}{lrrccc}
\toprule
\textbf{Label} & \textbf{TTK} & \textbf{Cyber} & \textbf{$\Delta F_1$} & \textbf{$\Delta P$} & \textbf{$\Delta R$} \\
\midrule
Émo & 5\,374 & 399 & \cellcolor[HTML]{C7EA97}\textcolor{black}{+0.27} & \cellcolor[HTML]{CFEB8F}\textcolor{black}{+0.31} & \cellcolor[HTML]{C0EAA1}\textcolor{black}{+0.23} \\
Base & 4\,133 & 364 & \cellcolor[HTML]{E6EC7E}\textcolor{black}{+0.37} & \cellcolor[HTML]{CBEA92}\textcolor{black}{+0.29} & \cellcolor[HTML]{EEDA6E}\textcolor{black}{+0.43} \\
Montrée & 971 & 301 & \cellcolor[HTML]{EEE172}\textcolor{black}{+0.42} & \cellcolor[HTML]{EEDF71}\textcolor{black}{+0.42} & \cellcolor[HTML]{EEE373}\textcolor{black}{+0.41} \\
Colère & 1\,180 & 290 & \cellcolor[HTML]{EFCD67}\textcolor{black}{+0.46} & \cellcolor[HTML]{EEDF71}\textcolor{black}{+0.42} & \cellcolor[HTML]{F1BB5E}\textcolor{black}{+0.49} \\
Dégoût & 52 & 80 & \cellcolor[HTML]{BBEAB5}\textcolor{black}{+0.14} & \cellcolor[HTML]{B62020}\textcolor{white}{+1.00} & \cellcolor[HTML]{C3ECC9}\textcolor{black}{+0.08} \\
Autre & 1\,270 & 62 & \cellcolor[HTML]{E4602C}\textcolor{white}{+0.76} & \cellcolor[HTML]{E25727}\textcolor{white}{+0.78} & \cellcolor[HTML]{E97B39}\textcolor{black}{+0.68} \\
Désignée & 1\,499 & 55 & \cellcolor[HTML]{EFA24F}\textcolor{black}{+0.56} & \cellcolor[HTML]{F1B55B}\textcolor{black}{+0.50} & \cellcolor[HTML]{ED9447}\textcolor{black}{+0.60} \\
Suggérée & 1\,909 & 54 & \cellcolor[HTML]{E5682F}\textcolor{black}{+0.73} & \cellcolor[HTML]{E97C3A}\textcolor{black}{+0.67} & \cellcolor[HTML]{E5642D}\textcolor{black}{+0.74} \\
Comportementale & 1\,242 & 36 & \cellcolor[HTML]{E77134}\textcolor{black}{+0.71} & \cellcolor[HTML]{EC9044}\textcolor{black}{+0.62} & \cellcolor[HTML]{E4602C}\textcolor{white}{+0.76} \\
Joie & 888 & 30 & \cellcolor[HTML]{F1BA5D}\textcolor{black}{+0.49} & \cellcolor[HTML]{F1B95D}\textcolor{black}{+0.49} & \cellcolor[HTML]{F1BB5E}\textcolor{black}{+0.49} \\
Complexe & 568 & 14 & \cellcolor[HTML]{EFA451}\textcolor{black}{+0.55} & \cellcolor[HTML]{E87637}\textcolor{black}{+0.69} & \cellcolor[HTML]{CDEB90}\textcolor{black}{+0.30} \\
Tristesse & 673 & 14 & \cellcolor[HTML]{E87737}\textcolor{black}{+0.69} & \cellcolor[HTML]{E4622D}\textcolor{black}{+0.75} & \cellcolor[HTML]{ED9347}\textcolor{black}{+0.60} \\
Jalousie & 7 & 6 & \cellcolor[HTML]{D3EEDE}\textcolor{black}{0.00} & \cellcolor[HTML]{D3EEDE}\textcolor{black}{0.00} & \cellcolor[HTML]{D3EEDE}\textcolor{black}{0.00} \\
Surprise & 824 & 5 & \cellcolor[HTML]{F0A853}\textcolor{black}{+0.54} & \cellcolor[HTML]{EE9B4B}\textcolor{black}{+0.58} & \cellcolor[HTML]{F1BB5E}\textcolor{black}{+0.49} \\
Peur & 1\,047 & 4 & \cellcolor[HTML]{EA833D}\textcolor{black}{+0.66} & \cellcolor[HTML]{E5662E}\textcolor{black}{+0.74} & \cellcolor[HTML]{EDED79}\textcolor{black}{+0.39} \\
Culpabilité & 19 & 3 & \cellcolor[HTML]{D3EEDE}\textcolor{black}{0.00} & \cellcolor[HTML]{D3EEDE}\textcolor{black}{0.00} & \cellcolor[HTML]{D3EEDE}\textcolor{black}{0.00} \\
Fierté & 202 & 3 & \cellcolor[HTML]{C9ECD1}\textcolor{black}{-0.05} & \cellcolor[HTML]{C0EBC4}\textcolor{black}{+0.09} & \cellcolor[HTML]{BFEAA3}\textcolor{black}{-0.22} \\
Embarras & 162 & 2 & \cellcolor[HTML]{E9803C}\textcolor{black}{+0.66} & \cellcolor[HTML]{E4632D}\textcolor{black}{+0.75} & \cellcolor[HTML]{BFEAA3}\textcolor{black}{-0.22} \\
Admiration & 211 & 0 & \cellcolor[HTML]{E5632D}\textcolor{black}{+0.75} & \cellcolor[HTML]{D93F20}\textcolor{white}{+0.85} & \cellcolor[HTML]{E9803C}\textcolor{black}{+0.66} \\
\midrule
\rowcolor[HTML]{3B3B3F}
\multicolumn{6}{l}{\textcolor{white}{\textbf{Métriques globales}}} \\
Macro F1 & 27\,911 & 781 & \cellcolor[HTML]{EFCC66}\textcolor{black}{+0.4570} & & \\
Micro F1 & 27\,911 & 781 & \cellcolor[HTML]{EEDC6F}\textcolor{black}{+0.4246} & & \\
\bottomrule
\end{tabular}
%
  }%
}

\newcommand{\metricbox}[3]{%
  \begin{minipage}[t]{0.31\textwidth}
    \setlength{\fboxsep}{0pt}%
    \colorbox{softSurface}{%
      \begin{minipage}[t][1.65cm][c]{\linewidth}
        \centering
        {\Large\bfseries #1}\par
        \vspace{0.15em}
        {\scriptsize\color{emotycMuted}#2}\par
        \vspace{0.15em}
        {\scriptsize #3}
      \end{minipage}%
    }
  \end{minipage}%
}

\title{Transferabilit\'e d'EMOTYC}
\subtitle{Lecture de la heatmap des deltas TextToKids $\rightarrow$ CyberAggAdo}
\author{}
\date{}

\begin{document}

\begin{frame}
  \titlepage
  \vspace{-2.2em}
  \begin{center}
    \small
    Heatmap source : \texttt{results/heatmap\_delta.html}\\
    Script source : \texttt{visualizations/delta\_heatmap.py}
  \end{center}
\end{frame}

\begin{frame}{Ce que mesure la heatmap}
  \begin{columns}[T,onlytextwidth]
    \begin{column}{0.58\textwidth}
      \[
        \Delta_m(\ell)=m_{\text{TextToKids}}(\ell)-m_{\text{CyberAggAdo}}(\ell)
      \]

      \begin{itemize}
        \item Les colonnes comparent $F_1$, la pr\'ecision et le rappel, label par label.
        \item La couleur encode $|\Delta|$ : vert clair si les performances sont proches, orange/rouge si l'\'ecart est fort.
        \item Une valeur positive signifie que le mod\`ele performe mieux sur TextToKids que sur CyberAggAdo.
      \end{itemize}
    \end{column}
    \begin{column}{0.36\textwidth}
      \vspace{0.4em}
      \begin{tabular}{@{}lr@{}}
        \toprule
        Corpus & Unit\'es textuelles \\
        \midrule
        \textcolor{ttkBlue}{TextToKids} & 27\,911 \\
        \textcolor{cyberOrange}{CyberAggAdo} & 781 \\
        \bottomrule
      \end{tabular}

      \vspace{1.2em}
      \small
      Configuration : template avec contexte, seuil 0.5, comparaison directe des r\'esum\'es JSON.
    \end{column}
  \end{columns}
\end{frame}

\begin{frame}{Heatmap compl\`ete des deltas}
  \vspace{-0.7em}
  \begin{center}
    \scriptsize
    \setlength{\tabcolsep}{4pt}
    \renewcommand{\arraystretch}{0.82}
    \resizebox{0.96\textwidth}{!}{\InputHeatmapTable}
  \end{center}

  \vspace{-0.4em}
  \scriptsize
  \textcolor{emotycMuted}{Lecture : $\Delta>0$ indique une performance sup\'erieure sur TextToKids ; la couleur indique l'amplitude, pas le sens.}
\end{frame}

\begin{frame}{Signal principal : forte perte hors domaine}
  \metricbox{$+0.4570$}{Macro $F_1$}{\'ecart TextToKids -- CyberAggAdo}
  \hfill
  \metricbox{$+0.4246$}{Micro $F_1$}{\'ecart TextToKids -- CyberAggAdo}
  \hfill
  \metricbox{781}{unit\'es CyberAggAdo}{contre 27\,911 pour TextToKids}

  \vspace{1.5em}
  \begin{itemize}
    \item Le mod\`ele conserve une d\'etection globale de l'\'emotion, mais les performances baissent nettement en transfert.
    \item Les labels fr\'equents dans CyberAggAdo montrent d\'ej\`a un d\'ecrochage substantiel : \textit{\'Emo} ($\Delta F_1=+0.27$), \textit{Base} ($+0.37$), \textit{Montr\'ee} ($+0.42$), \textit{Col\`ere} ($+0.46$).
    \item Le d\'ecalage touche \`a la fois la pr\'ecision et le rappel, ce qui sugg\`ere des faux positifs et des faux n\'egatifs hors domaine.
  \end{itemize}
\end{frame}

\begin{frame}{Labels qui concentrent le d\'ecrochage}
  \small
  \begin{tabular}{@{}lrrrrr@{}}
    \toprule
    Label & Cyber & TTK & $\Delta F_1$ & $\Delta P$ & $\Delta R$ \\
    \midrule
    Autre & 62 & 1\,270 & +0.76 & +0.78 & +0.68 \\
    Sugg\'er\'ee & 54 & 1\,909 & +0.73 & +0.67 & +0.74 \\
    D\'esign\'ee & 55 & 1\,499 & +0.56 & +0.50 & +0.60 \\
    Col\`ere & 290 & 1\,180 & +0.46 & +0.42 & +0.49 \\
    Montr\'ee & 301 & 971 & +0.42 & +0.42 & +0.41 \\
    Base & 364 & 4\,133 & +0.37 & +0.29 & +0.43 \\
    \bottomrule
  \end{tabular}

  \vspace{1.1em}
  \begin{itemize}
    \item Le mode \textit{Sugg\'er\'ee} d\'ecroche fortement : les indices contextuels appris sur TextToKids transf\`erent mal vers les messages de cyberharc\`element.
    \item \textit{Autre} est le plus grand \'ecart parmi les labels avec au moins 50 occurrences positives dans CyberAggAdo.
    \item \textit{Col\`ere} et \textit{Montr\'ee} restent suffisamment repr\'esent\'es pour signaler une vraie difficult\'e, pas seulement un bruit d'effectif.
  \end{itemize}
\end{frame}

\begin{frame}{Points de vigilance}
  \begin{itemize}
    \item Les labels tr\`es rares dans CyberAggAdo doivent \^etre interpr\'et\'es avec prudence : \textit{Admiration} a 0 occurrence positive, \textit{Embarras} 2, \textit{Fiert\'e} 3, \textit{Peur} 4.
    \item La couleur ne distingue pas le signe : une case claire peut aussi correspondre \`a une absence de cas suffisants pour comparer.
    \item La heatmap sert surtout \`a localiser les zones de non-transf\'erabilit\'e ; l'analyse qualitative des erreurs reste n\'ecessaire pour expliquer les m\'ecanismes.
  \end{itemize}

  \vspace{1em}
  \small
  Conclusion : EMOTYC apprend bien le domaine TextToKids, mais les formes discursives de CyberAggAdo changent assez pour d\'egrader la g\'en\'eralisation, notamment sur les modes d'expression et certains labels \'emotionnels.
\end{frame}

\end{document}
